
\subsubsection{2.1 Efficient Attention Mechanisms}

Standard transformer attention computes pairwise interactions across all positions, yielding O($n^{2})$ complexity \citep{vaswani2017attention}. Linear attention methods replace softmax with kernel approximations, achieving O(n) complexity (Katharopoulos et al., 2020; Choromanski et al., 2021). Sparse attention restricts the attention pattern to predefined or learned subsets: Longformer combines local windowed attention with global tokens \citep{beltagy2020longformer}, and BigBird adds random attention for theoretical expressiveness guarantees \citep{zaheer2020big}. Low-rank attention approximates attention matrices via factorization \citep{wang2020linformer}.

\subsubsection{2.2 Iterative and Recurrent Attention}

Recurrent attention uses RNN-like dynamics within attention computation \citep{graves2016adaptive}. Universal Transformers apply transformer blocks iteratively with shared parameters \citep{dehghani2019universal}. Equilibrium models solve for fixed points in attention computation (Bai et al., 2019; Bai et al., 2020). These works propose that iterative processing allows refinement toward task-relevant patterns but do not test whether the proposed iterative refinement mechanism actually operates.

\subsubsection{2.3 Mechanism Validation in Deep Learning}

Interpretability research seeks to understand what networks learn (Olah et al., 2020; Elhage et al., 2021). Ablation studies test component contributions but rarely decompose claims into orthogonal sub-hypotheses with falsification criteria. Negative results are underreported in machine learning (Henderson et al., 2018; Dodge et al., 2019).

